Peter was there when Jesus told the story about the man going away. He and three other disciples had asked Jesus a question. They could see the Temple across from where they were sitting on the Mount of Olives. Jesus had spoken of the day when those stones would be thrown down. The disciples wanted to know when it would happen. Jesus ends his answer by asking them to imagine what might happen in a house where nobody knows when the master will come home.

The master puts his servants in charge before he leaves. Each has work to do, and the man at the door has a particular reason to stay awake. The others depend on him to hear their master when he comes home. Jesus names the hours at which the traveler might arrive. A servant who had waited faithfully until the cock crowed would still have to watch until morning if the master had not come. Then Jesus says that what he has said to these few disciples is meant for everyone. Mark tells us what happened when Jesus asked Peter and two of the others to watch with him.

In the garden of Gethsemane, Jesus takes Peter, James and John a little farther with him than the rest. He tells them how deeply troubled he is. There is something he wants from his friends while he goes on alone to pray. He asks them to stay there and keep watch. Then he goes a little farther and falls to the ground to pray to his Father. Only a short time before, Peter had told Jesus he was prepared to die with him. When Jesus comes back, he finds Peter asleep.

Jesus calls him Simon and asks about the one hour he was to have watched. Peter had been so sure of what he would do even if everyone else deserted Jesus. Here was something his friend had actually asked of him, and already the hour was gone. Jesus gives him something he can still do. He tells the disciples to pray for help against temptation. He knows their willingness. He knows, too, how weary they are. When he comes back again and finds them asleep, Mark tells us that they have no answer for him.

Isaiah gives us words for a silence like that. The people would like God to find them doing what is right. When they look at the lives they have been leading, they can hardly remember anyone taking the trouble even to call on him. They describe a leaf that has come loose. It goes wherever the wind takes it. The wind in Isaiah's prayer is their own guilt. They have done things they can remember choosing to do. Now those choices have such a hold on them that they feel themselves being carried along. Even when they try to point to something good they have done, they can see the harm that has got mixed up in it.

They call God their Father. They called him that before they began telling him what has happened to them, and they can still call him that when they have finished. At the end they speak about clay in a potter's hands. They belong to the one who has shaped them. Saint Jerome understood what gives them the courage to speak. Clay has no say in how a potter shapes it. But these people are also children, and a child in trouble can go to his father. They can tell him the truth about themselves, even the truth that they have got out of the habit of speaking to him. And already, in the very act of telling him that, they have begun to pray again.

In the psalm, Israel asks God to come and see a vine. His own hand put it in the ground. The people asking him to look after it are the very people he once brought out of Egypt and planted in a land of their own. Toward the end of the psalm they say that they will call on him again. Saint Augustine noticed which comes first. First they ask God to give them life. Then they will have the strength to call his name. Even the desire that stirs in someone who has been away from God a long time can be received with gratitude. God is already doing for that person something he did for Israel.

At the beginning of Mass the priest spoke to the Father about the day when we will meet Christ. We hope to come before him having done some good for people he loves. Even before we spoke of the deeds we hope to bring him, we asked God to give us the will to do them. There are times when someone has asked for our company and we would sincerely like to be the kind of person who would go. Then the time comes to leave the house. We know how long we may have to sit there and how little we can think of to say. There is a distance to go between loving the person who asked and being willing to give up that hour.

Jesus knew that distance in the friends he took with him into the garden. He went on to give his life for those very men. When we come to Communion, we receive the Lord who came back to his sleeping friends. The people beside us can hear our Amen. God knows what it has taken for each of us to say it. Some of us have said it with a heart full of gratitude. Some of us can barely bring ourselves to ask him to help us begin again. What we have received from him here can touch even our reluctance to see someone who is suffering.

If you have promised to sit with someone who is afraid to be alone, keep that visit, even if you still have no words for what the person is going through. In the garden Jesus himself wanted the company of friends who could do little for him except stay. He knows what it costs you to go, and he knows the comfort it can be to someone to see you come through the door.

At the beginning of Advent, you can come to the Father who made you with the truth about the hours you have lost. The Son who gave himself for you can bring to life in you the very love you wished you had when you made that promise. When he comes in glory, you will be looking into the face of someone who has known you in your weakness and loved you there.
